\cotefile{4}
\scanpage{2}
% letterhead
\begin{center}
INSTITUT\\
DES HAUTES ÉTUDES SCIENTIFIQUES\\
\rule{1.1cm}{.35pt}
\end{center}

8.4.4.

\uline{NB}

a) $F$ opérant sur $H^i(X)\otimes_W k$ pour $i>n$ est \uline{nilpotent}, donc aussi opérant sur $\mathrm{Tors}^i(X)\otimes_W k$ ou sur $\mathrm{Tor}^{W}_{1}(\deleted{\unsure{\mathrm{Tor}}}H^i(X),k)\simeq {}_p\mathrm{Tors}^i(X)$

On a $H^n(X,\underline{O}_X)_{ss}\simeq(H^n(X)\otimes_W k)_{ss}$ d'où
\[
\deleted{\unsure{(H^{n}_{\rm lib}(X)\otimes k)_{ss}\simeq{}}}
0\longrightarrow(\mathrm{Tors}^{n}_{\deleted{\unsure{n}}}(X)\otimes k)_{ss}
\longrightarrow H^n(X,\underline{O}_X)_{ss}
\unsure{\longrightarrow|-}(H^n_{\rm lib}(X)\deleted{{}_{ss}}\otimes k)_{ss}\longrightarrow0
\tag*{$(*)$}
\]

b) $V$ opérant sur $H^i(X)\otimes_W k$ pour $i<n$ est nilpotent, et aussi opérant sur $\mathrm{Tors}^i(X)\otimes_W k$ et $ {}_p\mathrm{Tors}^i(X)$ pour $\uuline{i\leq n}$. On a
\[
0\longrightarrow(H^n_{\rm lib}(X)\otimes k)_{\deleted{ss}}^{ss}
\longrightarrow H^0(X,\underline{\Omega}^{n}_X)^{ss}
\unsure{\Longrightarrow|\longrightarrow} {}_p\mathrm{Tors}^{n+1}(X)_{\deleted{ss}}^{ss}
\longrightarrow\deleted{\unsure{{}_p\mathrm{Tor}}}0
\tag*{$(**)$}
\]

Sans doute $(*)$ et $(**)$ sont transposés l'un de l'autre\unsure{...|---}

\uline{Notation} L'indice $ss$ indique partie semi-simple pour $F$

l'exposant $^{ss}$ \rule[.5ex]{.38\linewidth}{.3pt} $V$
% archive sheet number: 1

\scanpage{3}
8. \uline{Une \deleted{\unsure{esti}}estimation} (\uline{heuristique}) \uline{du nombre des valeurs propres de frobenius dans la cohomologie évanescente qui sont des unités $p$-adiques.}

8.1. Soit $X$ un schéma lisse et propre sur le corps parfait $k$ \inserted{de car. $p\handwritten{>}0$}, et soit $W$ l'anneau des vecteurs de Witt sur $k$. Si on dispose d'une bonne théorie \unsure{de} la cohomologie cristalline $H(X)$ (au sens des catégories dérivées), on devra avoir la relation suivante avec la cohomologie de De Rham
\[
\handwritten{H_{\rm DR}}(X)=\underset{\handwritten{\smile}}{\overset{\handwritten{\frown}}{H\,(}}X)\overset{\rm L}{\otimes}_W k
\tag*{(8.1.1)}
\]
(isomorphisme dans la catégorie dérivée des $k$-modules). On en déduit \unsure{une} formule des coéfficients universels: suites exactes
\[
0\handwritten{\longrightarrow} H^i(X)\otimes_W k\handwritten{\longrightarrow} H^i_{\rm DR}(X)\handwritten{\longrightarrow}\mathrm{Tor}^W_1(H^{i+1}(X),k)\handwritten{\longrightarrow}0\;.
\tag*{(8.1.2)}
\]

Ce sont des suites exactes fonctorielles en $X$, tout comme l'isomorphisme (8.1.1). Sauf erreur, \deleted{\illegible{1}} ceci marche dès maintenant pour $p\handwritten{\ne}2$, grâce à Berthelot.
\inserted{\handwritten{[ On conclut, grâce à 8.4.3,}}
\[
\inserted{\handwritten{0\longrightarrow(H^i(X)\otimes_W k)_{ss}\longrightarrow
\underset{\text{(nul si }i>n\text{)}}{H^i(X,\underline O_X)_{ss}}\longrightarrow
\underset{\text{nul si }i\geq n}{({}_p\mathrm{Tors}^{i+1}(\deleted{\unsure{k}}X))_{ss}}\longrightarrow0}}
\]
% margin note written sideways
\margin{\handwritten{et $0\longrightarrow(H^{\unsure{i+n|i}}(X)\otimes_W k)^{ss}\longrightarrow H^{\unsure{i|i-n}}(X,\underline\Omega_X^n)^{ss}\longrightarrow({}_p\mathrm{Tors}^{\unsure{i+n+1|i-n+1}}(X))^{ss}\longrightarrow0$ où l'indice $ss$ signifie la partie semi-simple pour $F_X=F_X^*\sigma$ et \deleted{où} l'exposant $ss$ \rule[.5ex]{2em}{.3pt} $V_X=\sigma^{-1}F_{X*}$ ]}}

Si maintenant $Y$ est une section hyperplane de \deleted{degré assez élevé de} $X$ (supposée maintenant $\handwritten{\subset}\mathrm P^r_k$), on trouve une suite exacte analogue pour $Y$, et un homomorphisme \inserted{\handwritten{$u=(u_0,u_1,u_2)$}} de la première dans la seconde. Le diagramme du serpent nous donne alors une suite exacte
\[
\inserted{\handwritten{0\longrightarrow\mathrm{Ker}\,u_0\longrightarrow\mathrm{Ker}\,u_1\longrightarrow}}
\handwritten{\mathrm{Ker}\,u_2\longrightarrow}E(Y)\otimes_W k\handwritten{\longrightarrow}
\underset{\handwritten{\substack{\Vert\\\mathrm{Coker}\,u_1}}}{E_{\rm DR}(Y)}
\handwritten{\longrightarrow}\mathrm{Coker}\,u_{\handwritten{2}}\handwritten{\longrightarrow}0\;,
\tag*{(8.1.3)}
\]
où on a posé
\[
E(Y)=\mathrm{Coker}(H^{n-1}(X)\xrightarrow{\handwritten{u_0}}H^{n-1}(Y))\;,\quad
E_{\rm DR}(Y)=\mathrm{Coker}(H_{\rm DR}^{n-1}(X)\xrightarrow{\handwritten{u_1}}H_{\rm DR}^{n-1}(Y))\;,
\]
$X$ étant supposée connexe de dimension $n$, enfin $u_{\handwritten{2}}$ désigne l'application
\[
u_{\handwritten{2}}\mathbin{:}\mathrm{Tor}^{\handwritten{W}}_1(H^n(X),k)\handwritten{\longrightarrow}\mathrm{Tor}^W_1(H^n(Y),k)\;.
\tag*{(8.1.4)}
\]

\deleted{\handwritten{8.2.} Noter que $\mathrm{Tor}^{\deleted{\unsure{n}}W}_1(-,k)$ est isomorphe à $\mathrm{Hom}\deleted{\unsure{(}}_W(k,-)$ (élements de $-$ annulés par $pW$), et que $\deleted{\unsure{H}^{n}(}\quad\mathrm{Tors}^{n}(Y)$ ($=$\deleted{\unsure{(}}sous-module de torsion de $Y$) est dual (si on admet la \uline{dualité de Poincaré} habituelle pour la cohomologie cristalline, y compris les formules pour la torsion) de $\mathrm{Tors}^{n-1}(Y)$, et}
% slanted
\margin{\handwritten{il est \unsure{maladroit|malcommode} d'utiliser ici la dualité de Poincaré \deleted{\illegible{1}}}}
% archive sheet number: 2
\scanpage{4}
\deleted{par suite le deuxième membre de (8.1.4) est canoniquement dual à $\mathrm{Tors}^{n-1}(Y)\otimes_W k$. De même le premier membre est canoniquement \deleted{isomorphe à} \inserted{\handwritten{dual de}} $\mathrm{Tors}^{n+1}(X)\otimes_W k$, et enfin l'homomorphisme $u$ est transposé de l'homomorphisme}
\[
\deleted{u_{\handwritten{2}}'\mathbin{:}\mathrm{Tors}^{n-1}(Y)\otimes_W k\handwritten{\longrightarrow}\mathrm{Tors}^{n+1}(X)\otimes_W k}
\tag*{\deleted{(8.2.1)}}
\]
\deleted{induit par l'homomorphisme de Gysin}
\[
\deleted{H^{n-1}(Y)\handwritten{\longrightarrow}H^{n+1}(X)\;.}
\]

\deleted{Par suite, les termes $\mathrm{Ker}\,u_2$ et $\mathrm{Coker}\,u_2$ dans (8.1.3) peuvent s'interpréter \deleted{comme} respectivement comme \deleted{\unsure{Coker}} les duals de $\mathrm{Coker}\,u'_2$ et $\mathrm{Ker}\,u'_2$}

8.3. Considérons l'homomorphisme induit par l'inclusion \handwritten{$Y\longrightarrow X$,}
\[
u_{\deleted{\unsure{1}}0}\mathbin{:}H^{n-1}(X)\handwritten{\longrightarrow}H^{n-1}(Y)\;.
\]

\deleted{Sauf erreur,} \deleted{d}\handwritten{D}ans la situation analogue \deleted{sur le corps des complexes,} \inserted{\handwritten{en cohomologie $l$-adique ($l\ne p$)}} \uline{les \deleted{assertions} \inserted{\handwritten{théorème}} de Lefschetz} \deleted{dans son livre ``L'Analysis Situs ...''} \inserted{\handwritten{faible (pour les coefficients $Z/l^{\deleted{\unsure{m|n}}}Z$)}} impliquent que \uline{le conoyau} $E(Y)$ \uline{est sans torsion}, ce qui implique en particulier que $u_0$ induit un \uline{isomorphisme}
\[
\deleted{\unsure{H}}\handwritten{T}\mathrm{ors}^{n-1}(X)\handwritten{\overset{\sim}{\longrightarrow}}\mathrm{Tors}^{n-1}(Y)\;.
\tag*{(8.3.1)}
\]

Nous admettrons que ces relations restent valables dans le cas de la cohomologie cristalline. Alors l'homomorphisme (8.2.1) s'interprète comme l'homomorphisme
\[
u''_2\mathbin{:}\mathrm{Tors}^{n-1}(X)\otimes_W k\handwritten{\longrightarrow}\mathrm{Tors}^{n+1}(X)\otimes_W k
\tag*{(8.3.2)}
\]
induit par l'homomorphisme $H^{n-1}(X)\handwritten{\longrightarrow}H^{n+1}(X)$, cup-produit par la classe de polarisation de $X\handwritten{\subset}\mathrm P^r_k$. La chose importante pour nous est que ce morphisme (donc son Ker ) ne dépend plus du choix de $Y$, donc \deleted{aussi du} \deleted{\handwritten{\unsure{il en est de même}}} \inserted{\handwritten{le}} dernier terme de (8.1.3) \deleted{ne dépend} (qui est le dual de ce Ker) ne dépend plus de $Y$. Un autre point important est que, $E(Y)$ étant sans torsion, il a même rang que $E(Y)\otimes_W k$ intervenant dans (8.1.3) (et en particulier, si on a un endomorphisme de $E(Y)$, par exemple un frobenius, les valeurs propres de l'endomorphisme induit de $E(Y)\otimes k$ sont bien les réductions
% archive sheet number: 3
\scanpage{5}
mod $p$ des valeurspropres de l'endomorphisme donné de $E(Y)$ ). Enfin, (et c'est un petit luxe supplémentaire moins important pour nous, car seul les trois derniers termes de (8.1.3) sont pour nous vraiment importants du point de vue assymptotique ...) \deleted{\unsure{on aura} $\mathrm{Ker}\,u$} le \deleted{\unsure{premeir}}premier terme $\mathrm{Ker}\,u_0$ dans (8.1.3) est nul:
\[
\mathrm{Ker}\,u_0=0\;.
\tag*{(8.3.3)}
\]

8.4. Notons maintenant que l'on a des applications semi-linéaires $\mathrm{Frob}$ opérant fonctoriellement sur tous les espaces de cohomologie envisagés, et en particulier sur \deleted{\unsure{ceux}} \inserted{les termes} intervenant dans (8.1.3), qui sont donc des vectoriels de dimension finie sur $k$ avec endomorphisme $p$-linéaire. Considérons, pour un tel vectoriel $V$, le plus grand sous-espace vectoriel stable $V_{ss}$ (partie ``semi-simple'' de $V$) stable par $F_V$, et sur lequel $F_V$ soit \inserted{``semi-simple'' i.e.} surjectif (NB $k$ est parfait); \deleted{Le foncteur} il revient au,même de dire que $F$ est surjectif sur ce sous-espace, et \deleted{\unsure{nul sur le}}\uline{nilpotent} sur le quotient. Pour $V$ variable, la formation de $V_{ss}$ est un foncteur exact en $V$\unsure{.} Ainsi, la suite exacte (8.1.3) nous donne, compte tenu de (8.3.3), la suite exacte d'espace vectoriels à opérateurs $p$-linéaires semi-simples:
\[
\begin{aligned}
0&\handwritten{\longrightarrow}(\mathrm{Ker}\,u_1)_{ss}\handwritten{\longrightarrow}(\mathrm{Ker}\,u_2)_{ss}
\inserted{\handwritten{=({}_p(\mathrm{Tors}^n(X)))_{ss}}}\handwritten{\longrightarrow}\\
& (E(Y)\otimes_W k)_{ss}\handwritten{\longrightarrow}E_{\rm DR}(Y)_{ss}\handwritten{\longrightarrow}
\overset{\handwritten{\substack{0\\\Vert}}}{(\mathrm{Coker}\,u_2)_{ss}}\handwritten{\longrightarrow}0\;.
\end{aligned}
\tag*{(8.4.1)}
\]
\margin{\handwritten{en vertu de 8.4.4.}}

\inserted{\handwritten{8.4.2.}} D'après la théorie de structure des vectoriels à opération $p$-linéaire semi-simple, \deleted{qui} \inserted{ceux-ci} peuvent s'interpréter aussi comme des Modules libres de type fini sur le faisceau d'anneaux étale constant $\underline{\mathrm F}_p$ (corps premier à $p$ éléments), ou encore \deleted{\unsure{comme}} (une fois choisie une clôture séparable) comme des \deleted{\unsure{$G$-$\underline{\mathrm F}$}} représentations de $G=\mathrm{Gal}(\bar k/k)$ sur des vectoriels de dimension finie sur $\handwritten{\unsure{\mathrm F|\mathbb F|\underline{\mathrm F}}_p}$. Dans cette catégorie il y a une autodualité, qui s'interprète indffféremment en termes de dualité d'espaces de dimension finie sur $k$ avec opération $p$-linéaire semi-simple, ou en termes de dualité de
% archive sheet number: 4
\scanpage{6}
\deleted{\unsure{$G$-$\underline{\mathrm F}$ $G$-$\underline{\mathrm F}_p$}($G$)-Modules} Modules de représentations de $G$ sur $\underline{\mathrm F}_p$. Alors les termes $(\mathrm{Coker}\,u_2)_{ss}$ et $(\mathrm{Ker}\,u_2)_{ss}$ peuvent s'interpréter aussi comme les duals de $(\mathrm{Ker}\,u''_2)_{ss}$ et de $(\mathrm{Coker}\,u''_2)_{ss}$, $u''_2$ défini dans (8.3.2). Enfin \deleted{$(\mathrm{Ker}\,u_1$} $(\mathrm{Coker}\,u_1)_{ss}$ et $(\mathrm{Ker}\,u_1)_{ss}$ ont aussi une interprétation simple, grâce au

\uline{Lemme} 8.4.\deleted{\unsure{2}}\handwritten{3}. Soit $X$ un schéma propre sur un corps $k$, de car. $p\handwritten{>}0$. Alors pour tout $i\handwritten{\in}\underline Z$, l'homomorphisme d'augmentation \deleted{\unsure{$H$}}
\[
H^i_{\rm DR}(X)\handwritten{\longrightarrow}H^i(X,\underline O_X)
\]
induit un isomorphisme sur les parties semi-simples (pour l'opération $p$-linéaire de Frobenius).

Ceci résulte aussitôt de la suite spectrale de Hodge
\[
E^{a,b}_1=H^b(X,\underline\Omega^a_{X/k})\handwritten{\Longrightarrow}H^{\bullet}_{\rm DR}(X)\;,
\]
compte tenu que $\mathrm{Frob}$ opère \deleted{semi} $p$-linéairement sur cette suite spectrale \inserted{et} qu'il est nul sur $E_1^{a,b}$ pour $a\handwritten{\ne}0$. Alors les $(E_1^{0,b})\deleted{=}_{ss}=H^b(X,\underline O_X)_{ss}$ sont des cycles universels pour les $d_r$, et on gagne ...
% margin note written sideways
\margin{\handwritten{Cor. 8.4.\uline{$4\unsure{\mathrm{bis}|\mathrm{ter}}$} (utilisant 8.1.\unsure{?|2}) les résultats de la feuille jointe (marquée 8.4.4).}}

Mettant ensemble les renseignements obtenus, \deleted{et en même temps, pour des raisons de commodité, transformant par dualité de Poincaré l'homomorphisme $u''_2$ de (8.3.2), de sorte qu'on trouve (oubli dans 8.3 !)}
\[
\deleted{u_2\quad\text{est can. isomorphe à}}
\tag*{\deleted{(8.4.2)}}
\]
on trouve:

``\uline{Théorème}'' 8.5. En admettant les\handwritten{\,}ingrédients précédents pour la cohomologie cristalline (\inserted{formule des} coefficients universels, dualité de Poincaré, théor\deleted{\unsure{ème}}\handwritten{ème} de Lefschetz \handwritten{``faible''}\deleted{, y compris pour les groupes de torsion} \inserted{\handwritten{sans négliger la torsion}}) on obtient une suite exacte\handwritten{,} relative à une variété lisse projective connexe $X\handwritten{\subset}\mathrm P^r_k$ \inserted{de di\deleted{\unsure{p}}mension $n$} sur un corps parfait, de car. $p\handwritten{>}0$, et une section hyperplane (ou par une hypersurface \inserted{\handwritten{de degré $d\geq1$}}, cela s'y ramène ...) $Y$ lisse de dimension $n-1$:
\[
\begin{aligned}
&\deleted{(8.5.1)\quad 0\quad (\mathrm{Ker}\,\unsure{v})\quad(\unsure{\mathrm{Coker}}\,\unsure{w})\unsure{'}}\\
&\deleted{(E(Y)\otimes_W k)_{ss}\quad\deleted{\mathrm{Coker}\,v}\unsure{E_{\rm DR}}}
\end{aligned}
\]
% archive sheet number: 5

\scanpage{7}
\[
\begin{aligned}
0&\handwritten{\longrightarrow}\mathrm{Ker}\,v\xrightarrow{\handwritten{w}}\deleted{\unsure{(\mathrm{Coker}\,w)'}}\inserted{\handwritten{({}_p\mathrm{Tors}^{n}(X))_{ss}}}\handwritten{\longrightarrow}\\
&(E(Y)\otimes_W k)_{ss}\xrightarrow{\handwritten{\varphi}}
\deleted{\mathrm{Coker}\,v}\underset{\handwritten{\substack{\Vert\\\mathrm{Coker}\,v}}}{\handwritten{E(Y,\underline O_Y)_{ss}}}\handwritten{\longrightarrow}
\deleted{\unsure{(\mathrm{Ker}\,w)'\longrightarrow0}}\handwritten{0}\;,
\end{aligned}
\tag*{(8.5.1)}
\]
où on désigne par $v$ l'homomorphisme de restriction
\[
v\mathbin{:}H^{n-1}(X,\underline O_X)_{ss}\handwritten{\longrightarrow}H^{n-1}(Y,\underline O_Y)_{ss}\;,
\tag*{(8.5.2)}
\]
(en particulier $\mathrm{Coker}\,v$ est la ``cohomologie évanescente cohérente'' $E(Y,\underline O_Y)$), \deleted{et par $w$ l'homomorphisme}
\margin{\handwritten{\deleted{\unsure{$\xi=c_X$|$\xi=c_1X$}}}}
\[
\deleted{w\mathbin{:}(\mathrm{Tors}^{n-1}(X)\otimes_W k)_{ss}\xrightarrow{\handwritten{\unsure{\cup d\xi|u_{d\xi}}}}(\mathrm{Tors}^{n+1}(X)\otimes_W k)_{ss}}
\tag*{\deleted{(8.5.3)}}
\]
\margin{\handwritten{Katz fait \unsure{observer} que \unsure{$V$} \illegible{1} \unsure{est tjrs} \unsure{nul} \unsure{?} sur la \illegible{1}}}
\deleted{induit par l'homomorphisme $H^{n\deleted{\unsure{+1}}\handwritten{-1}}(X)\handwritten{\longrightarrow}H^{n+1}(X)$, \deleted{\unsure{cup-}}défini par le cup-produit avec la classe de polarisation \inserted{\handwritten{$d\xi$}} sur $X$.} Enfin, pour utiliser cette suite exacte (8.5.1), on tiendra compte du fait que
\[
E(Y)\overset{\rm dfn}{=}\mathrm{Coker}(H^{n-1}(X)\handwritten{\longrightarrow}H^{n-1}(Y))\quad\deleted{\text{est}}
\tag*{(8.5.4)}
\]
est un $W$-module libre \uline{sans torsion}; en particulier, si $k$ est un corps fini, le rang de $(E(Y)\otimes_W k)_{ss}$ est égal au nombre des valeurs propres de Frobenius géométrique relatif à $k$ opérant sur $\deleted{H}E^{n-1}(\deleted{\unsure{X}}Y)$ (la cohomologie cristalline évanescente de $Y$; ou ce qui revient au même, la cohomologie $\handwritten{l}$-adique pour n'importe quel $\handwritten{l\ne}p$ - sous réserve de l'algébricité de l'opération $\handwritten{\Lambda}$ sur lescohomologies cristallines et $\handwritten{l}$-adiques) \uline{qui sont des unités $p$-adiques}; les valeurs propres en question, réduites modulo $p$, sont égales aux valeurs\handwritten{\,}propres du frobenius $\deleted{\unsure{p}}\unsure{k|p}$-linéaire (relatif au corps de base fini $k$) opérant sur le terme $E(Y)\otimes_W k$ de (8.5.1), ou si on préfère, en termes de l'interprétation 8.4.2, aux valeurs propres d\deleted{u}e \inserted{l'inverse du} générateur de frobenius de $G=\mathrm{Gal}(\bar k/k)$ opérant sur le vectoriel de dimension finie sur $\underline{\rm F}_p$ correspondant.

Notons d'autre part l'égalité suivante déduite de (8.5.1):
\[
\begin{aligned}
\dim(E(Y)\otimes_W k)_{ss}={}&\dim E(Y,\underline O_Y)_{ss}\ \deleted{\unsure{-\dim\mathrm{Ker}\,w}}\\
&+(\dim\deleted{\mathrm{Coker}\,w}\inserted{\handwritten{({}_p\mathrm{Tors}^n(X))_{ss}}}-\dim\mathrm{Ker}\,v)\handwritten{,}\ \deleted{\unsure{+}}
\end{aligned}
\tag*{(8.5.5)}
\]
où le terme entre parenthèses du deuxième membre est $\handwritten{\geq}0$, de s\deleted{\unsure{o}}orte qu'\unsure{on}
% archive sheet number: 6
\scanpage{8}
a aussi l'inégalité suivante (résultant déjà de la suite exacte des trois derniers termes de (8.5.1)):
\[
\dim(E(Y)\otimes_W k)_{ss}\handwritten{\geq}\dim E(Y,\underline O_Y)_{ss}
\deleted{-\dim\mathrm{Ker}\,w}
\tag*{(8.5.6)}
\]
\[
\deleted{\dim E(Y,\underline O_Y)_{ss}-t^{n-1}(X,p)\;,}
\]
\deleted{où $t^i(X,p)$ désigne le nombre de \inserted{$p$-}coéfficients de torsion de $X$ (défini via la cohomologie cristalline de $X$ en dimension $i$).}

\uline{Remarques} 8.5.7. Je pense que l'homomorphisme $v$ de (8.5.2) est injectif\deleted{\unsure{,}} (ce qui simplifie en conséquence (8.5.1) et (8.5.5), puisque on aurait $\mathrm{Ker}\,v=0$ ); c'est vrai en tous cas si \inserted{\handwritten{$n=2$, ou si}} $Y$ est \deleted{de degré \unsure{assez|asse}} u\deleted{\unsure{j}}ne section par une hypersurface de degré $d$ assez grand pour qu'on ait
\[
H^{n-1}(X,\underline O_X(-d))=0\;.
\]

\deleted{D'autre part, on peut si on veut dans 8.5.5 \deleted{\unsure{bolquer}}bloquer ensemble les deux termes $\dim\mathrm{Coker}\,w-\dim\mathrm{Ker}\,w$, \deleted{car on aura comme} \inserted{pour trouver} la différence de rangs \deleted{$\dim\mathrm{Coker}\,w-\dim\mathrm{Ker}\,w$ \unsure{est}} des deux membres de (8.5.3).}

8.5.8. Lorsque \deleted{dans} l'on applique (8.5.5) en prenant pour $Y$ une section de $X$ avec une hypersurface de degré $d$ ``assez générale'' (de telle façon que le rang semi-simple de $E(Y,\underline O_Y)$ ait la valeur maxima), alors on trouve qu'au\deleted{\unsure{t}} terme $-\dim\mathrm{Ker}\,v$ près (sans doute toujours nul, cf. 8.5.\deleted{\unsure{6}}7) le deuxième membre est \uline{indépendant du choix} \deleted{de} \uline{particulier de \handwritten{$Y$}}\unsure{.} Cette remarque s'applique en particulier (8.5.7) si $d$ est assez grand. On en conclut que si $k$ est fini, \deleted{\unsure{les,}} on peut trouver un ouvert \inserted{$U_d$} de l'espace projectif des hypersurfaces de degré $d$, tel que pour \deleted{tout} \inserted{un} point \inserted{variable} de $U$ à coéfficients dans une extension finie $k'$ de $k$, le nombre de valeurs propres de Frobenius opérant sur la cohomologie évanescente de $Y_{\overline{k'}}$ et qui sont des unités $p$-adiques soit constant. \deleted{Mieux, quitte à agrandir un peut $k$, de sorte que les valeurs propres de $\mathrm{Frob}_k$} Nous désirons mon%
% archive sheet number: 7
\scanpage{9}%
rer que pour $d$ assez grand, ce nombre est non nul (ce qui impliquera \deleted{\unsure{al}} alors\deleted{\unsure{,}} - et est sans doute équivalent - que le motif $E(Y)$ pour tous ces $Y$ est de \uline{niveau} maximum $n-1$, ou encore que le motif $E(Y)$ pour $Y$ générique est de niveau $n-1$). Sous la forme de (8.5.6)\deleted{\unsure{,}}, on voit que ceci est assuré dès que l'on sait que
\[
\dim E(Y,\underline O_Y)_{ss}\handwritten{>0}\deleted{\unsure{\longrightarrow+\infty}}\qquad\text{pour}\qquad d\handwritten{\longrightarrow}+\handwritten{\infty}\;,
\]
où le dim est pris bien entendu avec sa valeur générique. \deleted{\handwritten{On sait seulement (cf. n\textsuperscript{o} 7) que $\dim E(Y,\underline O_Y)_{ss}>0$ pour $d$ grand, et que $\varphi\ne0$ pour $d$ grand (donc $E(Y)$ de niveau $n-1$) dès que \unsure{l'on satisfait à|l'on sait que} $\dim(\mathrm{Tors}^{n+1}(X)\otimes k)_{ss}<p-1$.}}
% slanted
\margin{\handwritten{Or cela est établi au n\textsuperscript{o} 7}}

8.5.9. Le corps de base $k$ étant à nouveau quelconque (même pas \unsure{nécessairement} parfait), désignons, pour un degré $d$ donné, par $U_d$ l'ouvert de l'espace projectif des hypersurfaces de degré $d$ dans $\mathrm P^r_k$, formé des hypersurfaces dont l'intersection avec $X$ est \deleted{non singulière} \inserted{lisse} de dimension $n-1$, et telles que $E(Y,\underline O_Y)_{ss}$ soit de rang maximum. Alors lessuites exactes (8.5.1), pour les $Y$ provenant de points de $U_d$, proviennent d'une suite exacte de Modules localement libres \inserted{sur $U_d$} à opération $p$-linéaire semi-simple, ou de façon équivalente (cf 8.4.2) à une suite exacte de systèmes locaux \inserted{de vectoriels sur $\underline{\mathrm F}_p$} sur $U_d$ pour la top\deleted{\unsure{l}}ologie étale, \deleted{\unsure{formée d}} ou enfin, une fois choisi un point géométrique de $U_d$ (par exemple via une clôture \unsure{séparable|ségarable} du point générique) une suite exacte de représentations \inserted{de $G=\handwritten{\pi}_1(U_d)$} dans des vectoriels de dimension finie sur $\underline{\mathrm F}_p$. \deleted{\unsure{On}} On trouve un homomorphisme \inserted{\handwritten{$\varphi$}} du système local défini par les $E(Y)\otimes_W k'$ \deleted{dans} \inserted{\handwritten{\uline{sur}}} celui défini par les $E(Y,\underline O_Y)$, dont le\deleted{s} noyau \deleted{et le conoyau sont des} \inserted{\handwritten{est un}} système\deleted{s} loca\deleted{ux}\handwritten{l} \uline{constant} (i.e. provenant du corps de base $k$). Lorsque $k$ est fini, on peut supposer, quitte à passer à une extension finie de $k$, que les valeurs propres du frobenius relatif à $k$ opérant sur ce\deleted{s} dernier\deleted{s} \deleted{\unsure{(aux}} sont toutes égales à 1 (il suffit que les valeurs propres de frobenius relatif à $k$ opérant sur \deleted{les} $H^{n\deleted{\unsure{-1}}}(X,\underline O_X)_{ss}$ \deleted{et sur les deux membres de (8.5.3)} soient
% archive sheet number: 8
\scanpage{10}%
toutes égales à 1\handwritten{)}. On trouve alors une description complète de la variation des valeurs propres mod $p$ de frobenius opérant sur $E(Y)$, pour $Y$ variable avec un point de $U_d$, en termes essentiellement de la seule représentation de $G$ dans un vectoriel de dimension finie sur $\underline{\mathrm F}_p$, correspondant aux $E(Y,\underline O_Y)_{ss}$ (donc essentiellement aux ``matrices de Hasse-Witt des $Y$ en dimension critique), - et bien entendu de la distribution des \deleted{frobenius} éléments de frobenius dans $G$ en termes d'un point variable de $U_d$.

8.5.10. L'homomorphisme \deleted{\unsure{$\varphi$}}fondamental $\varphi$ de (8.5.1) est \deleted{surjectif (resp.} injectif\deleted{)} \inserted{\handwritten{donc bijectif}} si et seulement si $w$ dans 8.5.\deleted{\unsure{3}}\handwritten{1} est \deleted{injectif (resp. si $w$ est} surjectif, \inserted{\handwritten{et il l'est si $({}_p\mathrm{Tors}^{n}(X))_{ss}=0$,}} le seulement si étant vrai sous la réserve $\mathrm{Ker}\,v=0$ de 8.5.7); il suffit en particulier pour ceci que $X$ n'ait pas de $p$-torsion en \deleted{caractéristique} dimension $n\deleted{\unsure{-1}}$ \deleted{(resp. en dimension \unsure{$n$}\deleted{\unsure{, ou}} $n+1$,} \handwritten{(}ou ce qui revient au même par dualité, en dimension $n\handwritten{+1}$\handwritten{)}. Lorsque $X$ se \deleted{Je pense que même si} relève (avec sa polarisation) en caractéristique nulle, ce\deleted{s}\handwritten{tte} condition\deleted{s} \deleted{s'exp} \inserted{suffisante\deleted{s}} doi\deleted{vent}\handwritten{t} aussi pouvoir s'exprimer (une fois qu'on disposera de fondements satisfaisants pour la cohomologie cristalline) par l\deleted{es}\handwritten{a} condition\deleted{s} analogue\deleted{s} sur \inserted{torsion\handwritten{ de}} la cohomologie $p$-\uline{adique} de la variété relevée (sur un corps algébriquement clos, ou si on veut, sur le corps des complexes ...): \deleted{on réécrit (8.5.3) en termes de $p$-torsion de cette variété tensorisée par $\underline Z/p\underline Z$ (mais en supprimant les indices $ss$ , qui n'ont plus de sens), et on demande que $w_p$ ainsi défini soit surjectif resp. injectif ....}

\deleted{Je pense que même si $n=\dim X=2$, il peut arriver que $\varphi$ ne soit pas surjectif, \deleted{\unsure{i.e}}et qu'on ne peut pas dans l'inégalité (8.5.6) supprimer le terme correctif $\dim\mathrm{Ker}\,w$ ou $t^{n-1}(X,p)$. En effet\deleted{\unsure{,}}, si $X$ est une}
% archive sheet number: 9

\scanpage{11}%
\deleted{Noter d'autre part qu'on a équivalence entre les conditions suivantes:}

\deleted{a) L'homomorphisme $\varphi$ est surjectif (resp. injectif) pour tout $d$ et tout point $s\in U_d$ (définissant un $Y$ correspondant).}

\deleted{b) Même condition, mais avec \uline{un} $d$ fixé multiple de $p$, et tel que \deleted{\unsure{c)}} l'on ait $\mathrm{Ker}\,v=0$ (condition sans doute automatiquement \unsure{vérifié|vérifiée} (8.5.7)).}

\deleted{c) $(\mathrm{Tors}^{n-1}(X)\otimes_W k)_{ss}=0$ (resp. $(\mathrm{Tors}^{\handwritten{n+1}}(X)\otimes_W k)_{ss}=0$ ).}

\inserted{8.5.\deleted{11}\handwritten{\unsure{10|11}}.} Il serait intéressant de donner des exemples où \deleted{ces conditions ne sont \unsure{pas} satisfaites. Noter que la condition non respée est toujours satisfaite si $n=2$ ($X$ une surface)? Dans le cas général,} \inserted{\handwritten{$\varphi$ n'est pas bijectif}} on a l'inégalité
\[
\dim_k H^{n-1}(X,\underline O)_{ss}=\dim_{\underline{\mathrm F}_p}H^{n-1}(\overline X,\underline Z/p\underline Z)\handwritten{\geq}\dim\bigl((H^{\handwritten{n-1}}(X)/\mathrm{Tors}^{n-1}(X))\otimes k\bigr)_{ss}
\tag*{\inserted{(8.5.1\deleted{1}0.1)}}
\]
(grâce à 8.1.2 et 8.4.3 ; en fait cela marche en toute dimension $i$, pas seulement $i=n-1$), et dire que $\varphi$ est bijectif signifie que l'inégalité précédente est une égalité. L'exemple de Serre au symposium de Mexico fournit un $X$ (quotient d'une intersection complète par un groupe $G=\underline Z/p\underline Z$ opérant sans points fixes) pour lequel on a $H^1(X)=0$ , mais $H^1(X,\underline Z/pZ)\unsure{\simeq|\doteqdot|=|\ne}\underline Z/\unsure{p\underline Z}$. Quitte à couper avec des hpyperplans, on trouve le même exemple avec $\dim X=2$ . Pour cet exemple, \deleted{(8.5.11.1) n'est pas une égalité, et par suite,} d'après \deleted{ce qui précède,} \inserted{8.1.2, on aura} $(\handwritten{{}_p}\mathrm{Tors}^{\handwritten{2}}(X)\deleted{\unsure{\otimes_W k}})_{ss}\simeq\overset{\handwritten{H^1(X,\underline O_{\unsure{X|Y}})_{ss}}}{\overset{\handwritten{\shortparallel}}{H^1(X,\underline Z/p\underline Z)\otimes k}}=\mathrm{Ker}\,\handwritten{\varphi}$ (quel que soit $d$ et la section $Y$ de $X$ de degré $d$, permettant de définir $\varphi$ ) , qui est non nul, donc $\varphi$ n'est pas injectif. \deleted{Si maintenant $Z$ est un schéma lisse projectif sur $k$ de dimension $n-2$ ($n$ donné $\geq\deleted{\unsure{5|3}}\handwritten{4}$), à cohomologie cristalline sans torsion \inserted{et $(H^{n-4}(Z)\otimes k)_{ss}\ne0$} (par exemple $\mathrm P^{n-2}_k$ \inserted{si $n$ pair ...}), alors, posant $X'=X\times Z$ , on voit (sous réserve de la \uline{formule de Kunneth} en cohomologie cristalline) que $(\mathrm{Tors}^{n-1}(X')\otimes k)_{ss}$\deleted{\unsure{ntient comme facteur direct}} est non nul, ce qui fournit des exemples où $\varphi$ n'est pas surjectif, pour tout $n\geq4$. Il y a sans doute aussi des exemples avec $n=3$, en cherchant bien ....}
% archive sheet number: 10

\scanpage{12}%
8.6. La suite exacte (8.5.9) de systèmes locaux sur $U_d$ correspondant aux suites exactes (8.5.1) met en évidence que le\deleted{s} systêmes loca\deleted{ux}\handwritten{l} \unsure{corres}\inserted{pondant \deleted{aux}\handwritten{au}} \deleted{\unsure{des}} $(E(Y)\otimes_W k)_{ss}$ \deleted{et aux $E(Y,\underline O_Y)_{ss}$ ont}\inserted{\handwritten{a}}, en plus d'un sous-système \deleted{resp. d'un système quotient} \inserted{\handwritten{du premier}} de nature essentiellement triviale (car provenant du corps de base, donc constant du point de vue géométrique) \deleted{\unsure{une}}\handwritten{\unsure{un}} \unsure{composant} \deleted{commun $\mathrm{Coim}\,\varphi=\mathrm{Im}\,\varphi$} \inserted{\handwritten{$E(Y,\underline O_Y)_{ss}$}}, dont on s'attend qu'il donne lieu à une représentation tout à fait non triviale de $G$ (NB si on n'aime pas les groupes fondamentaux généraux et les systèmes locaux, on peut tout interpréter en termes de la suite exacte 8.5.1 correspondant au \uline{point générique} $s$ de l'espace des hypersurfaces de degré $d$ dans $\mathrm P^r_k$ , et remplacer $G$ par le groupe $\mathrm{Gal}(\overline{k(s)}/k(s))$, - mais on perd alors les frobenius ...). On aimerait pouvoir préciser ce point. Par exemple, \inserted{\handwritten{Si $n\geq2$, cette représentation est-elle \uline{irréductible}, ou même}} l'image du groupe \unsure{fondamental} \inserted{\handwritten{$G$}} dans $\mathrm{Aut}(V)$ ($V$ étant le \deleted{\unsure{corps}} vectoriel de dimension finie sur $\underline{\mathrm F}_p$ où s'effectue la représentation envisagée de $G$) est-elle égale à $\mathrm{Aut}(V)$ tout entier ? \deleted{Si oui,} \inserted{\handwritten{(Ce n'est pas \deleted{\unsure{toujours vrai}} lorsque $X=\mathrm P_k^{r=1}$, cas qu'on peut traiter trivialement.) Lorsque la condition \unsure{ci-dessus} est remplie,}} alors dès que \inserted{\handwritten{\unsure{l'on a}}} $\dim V\geq2$ (ce qui est sans doute \deleted{\unsure{vérifié}}\handwritten{vrai} dès que $d$ est assez grand, $\dim V$ tendant vers $\infty$ avec $d$, cf. 8.5.8), tout élément de $\underline{\mathrm F}_p$ est trace d'un élément de $G$ opérant sur $V$. De ceci, \inserted{de la ``formule de Woodshole'' \deleted{\handwritten{\unsure{pour les points fixes}}}} et du théorème de Čebotareff, on déduirait alors, \deleted{\unsure{k}}\handwritten{(}si le corps $k$ est fini\unsure{)}, qu'on peut, pour n'importe quel ouvert non vide $U$ de l'espace des hypersurfaces de degré $d$ dans $\mathrm P^r_k$ ($d$ étant choisi comme dessus) trouver un point \deleted{\unsure{de}} $s$ de $U$ à coéfficients dans une extension finie $k'$ de $k$, définissant une section $Y$ telle que $\mathrm{card}\,Y(k')$ soit congru mod $p$ à un élément de $\underline{\mathrm F}_p$ donné à l'avance. Au moment d'écrire ces lignes, et malgré l'assistance de gens compétents comme Katz, il ne semble pas qu'un résultat de cette nature soit prouvé\deleted{,}\handwritten{(} même si $X=\mathrm P^{2\handwritten{=r}}$ i.e. pour des courbes planes \inserted{\handwritten{$Y$}} de grand degré\handwritten{)}. \deleted{Ce résultat serait déjà bien intéressant, car il impliquerait déjà une réponse affirmative à la conjecture du n\textsuperscript{o}6 sur le niveau de la cohomologie évanescente (dont il est question aussi dans 8.5.8).}
% archive sheet number: 11

